\section*{CHƯƠNG 6 - KẾT LUẬN}
\addcontentsline{toc}{section}{\numberline{}CHƯƠNG 6 -  KẾT LUẬN}
\setcounter{section}{6}
\setcounter{subsection}{0}
\setcounter{figure}{0}
\setcounter{table}{0}

%-------------Phần trình bày nội dung chương 6
Đề tài này đã thành công trong việc xây dựng một hệ thống tự động hóa quản lý và cấu hình mạng doanh nghiệp hoàn chỉnh, sử dụng ngôn ngữ lập trình Python kết hợp với thư viện Netmiko để tương tác với các thiết bị mạng được mô phỏng trên nền tảng GNS3-VM. Đây không chỉ là một bài tập học thuật đơn thuần, mà là một minh chứng rõ ràng cho thấy khả năng áp dụng công nghệ tự động hóa vào môi trường thực tế của doanh nghiệp hiện đại. Hệ thống đã được thiết kế với kiến trúc modular, cho phép dễ dàng bảo trì, mở rộng và tùy chỉnh theo nhu cầu cụ thể của từng tổ chức.

Về mặt kỹ thuật, đề tài đã triển khai thành công nhiều chức năng quan trọng trong quản lý mạng. Hệ thống có khả năng tự động hóa hoàn toàn quy trình cấu hình VLAN cho các Switch, bao gồm việc tạo VLAN, đặt tên và cấu hình các port ở chế độ trunk hoặc access. Đối với các Router, chương trình tự động gán địa chỉ IP cho các interface và subinterface, cấu hình encapsulation dot1q, và kích hoạt các interface một cách có hệ thống. Đặc biệt, việc thiết lập giao thức định tuyến OSPF được thực hiện tự động, bao gồm cấu hình router-id, quảng bá các network vào OSPF area, và kiểm tra trạng thái neighbor sau khi cấu hình. Tất cả các thao tác này được thực hiện thông qua việc gửi các lệnh CLI đến thiết bị qua giao thức SSH, nhờ vào khả năng của thư viện Netmiko trong việc mô phỏng tương tác con người với thiết bị mạng.

Một trong những điểm nổi bật của hệ thống là chức năng Backup và Restore thông minh. Chương trình có khả năng tự động sao lưu cấu hình của tất cả thiết bị với timestamp rõ ràng, giúp dễ dàng theo dõi lịch sử thay đổi cấu hình. Hệ thống lưu trữ file backup theo định dạng hostname\_backup\_YYYYMMDD\_HHMMSS.txt, cho phép quản lý và tra cứu nhanh chóng. Khi cần khôi phục, người dùng có thể chọn từ danh sách các file backup có sẵn, và hệ thống sẽ tự động làm sạch cấu hình bằng cách loại bỏ các dòng comment và metadata không cần thiết, sau đó áp dụng cấu hình theo từng batch để tránh tình trạng quá tải. Đặc biệt, chức năng mô phỏng lỗi và tự động khôi phục đã chứng minh tính ổn định và độ tin cậy của hệ thống, khi chương trình có thể phát hiện và khôi phục cấu hình một cách tự động mà không cần sự can thiệp của con người.

Hệ thống giám sát được tích hợp trong dự án cung cấp khả năng theo dõi trạng thái mạng một cách toàn diện. Chương trình có thể kiểm tra kết nối SSH đến tất cả thiết bị, đo thời gian phản hồi và kiểm tra uptime của từng thiết bị. Đối với interface, hệ thống tự động thu thập thông tin về trạng thái up/down, địa chỉ IP được gán và các thông số kỹ thuật khác. Trên các Router, chương trình giám sát bảng định tuyến, kiểm tra các route được học qua OSPF hoặc static route, và đánh giá tình trạng kết nối giữa các router trong mạng OSPF. Đối với Switch, hệ thống theo dõi cấu hình VLAN, trạng thái các trunk port và phân bổ VLAN trên từng access port. Tất cả thông tin giám sát này được ghi nhận vào file log chi tiết và tạo báo cáo dưới dạng CSV để phục vụ cho việc phân tích và báo cáo định kỳ.

GNS3-VM đóng vai trò là nền tảng mô phỏng mạnh mẽ trong dự án này, cho phép tạo ra một môi trường mạng doanh nghiệp phức tạp mà không cần đầu tư vào thiết bị phần cứng đắt tiền. Khác với các simulator đơn giản như Packet Tracer, GNS3 sử dụng các IOS image thực của Cisco, mang lại trải nghiệm gần giống với môi trường production. Điều này đặc biệt quan trọng vì nhiều tính năng và hành vi của thiết bị thực không được mô phỏng chính xác trên các simulator giản lược. GNS3-VM tích hợp tốt với Python thông qua giao thức SSH, cho phép các script tự động có thể kết nối và điều khiển thiết bị từ xa một cách linh hoạt. Môi trường ảo hóa này cũng giúp tiết kiệm đáng kể chi phí điện năng và không gian vật lý so với việc vận hành một lab với thiết bị thật, đồng thời cho phép dễ dàng snapshot và rollback khi cần thiết.
    